% Part: incompleteness
% Chapter: incompleteness-provability
% Section: first-incompleteness-thm

\documentclass[../../../include/open-logic-section]{subfiles}

\begin{document}

\olfileid[fr]{inc}{inp}{1in}

\olsection{Le premier théorème d'incomplétude}

Nous pouvons maintenant présenter la démonstration
originale du premier théorème d'incomplétude de
Gödel. Soit $\Th{T}$ une théorie axiomatisée par un
ensemble calculable, dans un langage qui étend
celui de l'arithmétique, et qui contient les axiomes
de~$\Th{Q}$. Elle représente donc, en particulier,
les fonctions totales et les relations calculables.

Nous avons expliqué que, pour un codage raisonnable
des !!{formula}s et des preuves par des nombres,
la relation $\Prf[\Th{T}](x,y)$ est calculable.
Elle signifie que $x$ est le numéro de Gödel d'une
!!{derivation}, dans~$\Th{T}$, de la !!{formula}
dont le numéro de Gödel est~$y$. Pour la théorie
particulière qu'il considérait, Gödel a même montré
que cette relation est récursive primitive, au moyen
des 45 fonctions et relations énumérées dans son
article. La quarante-cinquième, $x B y$, est
précisément $\Prf[\Th{T}](x,y)$ pour ce choix de
$\Th{T}$. Rappelons que là où Gödel écrivait
« récursif », nous dirions aujourd'hui « récursif
primitif ».

La relation $\Prf[\Th{T}](x,y)$ étant calculable,
elle est représentable dans~$\Th{T}$. Notons
$\OPrf[\Th{T}](x,y)$ une !!{formula} qui la
représente et posons
$\OProv[\Th{T}](y)=
\lexists[x][\OPrf[\Th{T}](x,y)]$.
Cette !!{formula} décrit la quarante-sixième
relation de la liste de Gödel, $\fn{Bew}(y)$.
Selon Gödel, c'est la seule relation dont on ne
« peut affirmer qu'elle est récursive ». Il voulait
probablement dire que sa définition ne permet pas
de conclure à sa calculabilité. Pour une théorie cohérente
contenant~$\Th{Q}$, sa relation de démontrabilité n'est effectivement
pas calculable.

Supposons maintenant que $\Th{T}$ soit une théorie
!!{axiomatizable} contenant~$\Th{Q}$. La relation
$\Prf[\Th{T}](x,y)$ est décidable; une
!!{formula}~$\OPrf[\Th{T}](x,y)$ la représente
donc dans~$\Th{Q}$. Gardons la notation
$\OProv[\Th{T}](y)$ ci-dessus. Le lemme du point
fixe fournit un !!{sentence}~$!G_\Th{T}$ tel
que $\Th{Q}$, donc aussi $\Th{T}$, démontre
\begin{equation}
\ollabel{eqn:qpf}
!G_\Th{T}\liff\lnot\OProv[\Th{T}](\gn{!G_\Th{T}}).
\end{equation}
L'!!{sentence}~$!G_\Th{T}$ affirme donc, en
substance : « $!G_\Th{T}$ n'est pas démontrable
dans~$\Th{T}$ ».

\begin{lem}
\ollabel{lem:cons-G-unprov}
Si $\Th{T}$ est une théorie cohérente,
!!{axiomatizable} et étendant~$\Th{Q}$, alors
$\Th{T}\Proves/ !G_\Th{T}$.
\end{lem}

\begin{proof}
Supposons que $\Th{T}$ démontre $!G_\Th{T}$.
Il existe alors effectivement une !!{derivation}
de cet !!{sentence}; pour un certain entier~$m$,
la relation
$\Prf[\Th{T}](m,\Gn{!G_\Th{T}})$ est vraie.
Par représentation, $\Th{Q}$ démontre
$\OPrf[\Th{T}](\num m,\gn{!G_\Th{T}})$,
donc aussi
$\lexists[x][\OPrf[\Th{T}](x,\gn{!G_\Th{T}})]$,
c'est-à-dire
$\OProv[\Th{T}](\gn{!G_\Th{T}})$.
L'équivalence ci-dessus (\olref{eqn:qpf}) permet
alors à $\Th{Q}$ de démontrer $\lnot !G_\Th{T}$;
$\Th{T}$, qui étend $\Th{Q}$, le démontre aussi.
Si $\Th{T}$ démontrait $!G_\Th{T}$, elle
démontrerait donc aussi sa négation, en contradiction
avec sa cohérence.
\end{proof}

\begin{defn}
\ollabel{thm:oconsis-q}
Une théorie~$\Th{T}$ est \emph{$\omega$-cohérente}
si, pour tout !!{sentence} de la forme
$\lexists[x][!A(x)]$, lorsqu'elle démontre
$\lnot !A(\num 0)$, $\lnot !A(\num 1)$,
$\lnot !A(\num 2)$, \dots, elle ne démontre pas
$\lexists[x][!A(x)]$.
\end{defn}

Toute théorie $\omega$-cohérente est cohérente.
En effet, une théorie incohérente démontre tout
!!{sentence}; elle démontre donc, pour chaque~$n$,
$\lnot !A(\num n)$, ainsi que
$\lexists[x][!A(x)]$. Elle est donc
$\omega$-incohérente. Par contraposition,
la $\omega$-cohérence implique la cohérence.

\begin{lem}
\ollabel{lem:omega-cons-G-unref}
Si $\Th{T}$ est une théorie $\omega$-cohérente,
!!{axiomatizable} et étendant~$\Th{Q}$, alors
$\Th{T}\Proves/\lnot !G_\Th{T}$.
\end{lem}

\begin{proof}
Montrons que, si $\Th{T}$ démontre
$\lnot !G_\Th{T}$, elle est $\omega$-incohérente.
Supposons donc cette négation démontrable.
Si $\Th{T}$ est incohérente, le résultat est
immédiat. Sinon, elle est cohérente; le lemme
précédent (\olref{lem:cons-G-unprov}) montre
qu'elle ne démontre pas $!G_\Th{T}$.
Il n'existe donc aucune !!{derivation} de cet
!!{sentence} dans~$\Th{T}$. Par la clause
négative de représentation, $\Th{Q}$ démontre
\[
\lnot\OPrf[\Th{T}](\num 0,\gn{!G_\Th{T}}),
\lnot\OPrf[\Th{T}](\num 1,\gn{!G_\Th{T}}),
\lnot\OPrf[\Th{T}](\num 2,\gn{!G_\Th{T}}),\dots
\]
et $\Th{T}$ les démontre aussi. D'autre part,
par l'équivalence ci-dessus (\olref{eqn:qpf}),
$\lnot !G_\Th{T}$ équivaut à
$\lexists[x][\OPrf[\Th{T}](x,\gn{!G_\Th{T}})]$.
La théorie $\Th{T}$ est donc $\omega$-incohérente.
\end{proof}

\begin{prob}
Toute théorie $\omega$-cohérente est cohérente.
Montrez que la réciproque est fausse : il existe
des théories cohérentes qui ne sont pas
$\omega$-cohérentes. Pour cela, montrez que
$\Th{Q}\cup\{\lnot !G_\Th{Q}\}$ est cohérente,
mais $\omega$-incohérente.
\end{prob}

\begin{thm}
\ollabel{thm:first-incompleteness}
Toute théorie~$\Th{T}$ $\omega$-cohérente,
!!{axiomatizable} et étendant~$\Th{Q}$ est
incomplète.
\end{thm}

\begin{proof}
Si $\Th{T}$ est $\omega$-cohérente, elle est
cohérente. Le premier lemme
(\olref{lem:cons-G-unprov}) donne donc
$\Th{T}\Proves/ !G_\Th{T}$. Le second
(\olref{lem:omega-cons-G-unref}) donne
$\Th{T}\Proves/\lnot !G_\Th{T}$.
La théorie est incomplète, car elle ne démontre
ni $!G_\Th{T}$ ni sa négation.
\end{proof}

\end{document}
